\section{One-Pass versus Multi-Pass dan Organisasi Front-End/Back-End}
\label{sec:ch03-pass}

\subsection{Kompilator Satu Lintasan}

Kompilator satu lintasan (\eng{one-pass compiler}) membaca program
sumber satu kali dan langsung menghasilkan kode target. Setiap
konstruksi bahasa harus dapat diproses tanpa menunggu informasi dari
bagian program yang belum dibaca. Pendekatan ini cocok untuk bahasa
yang dirancang agar dapat dikompilasi satu lintasan, seperti Pascal
klasik.

Keuntungan one-pass adalah kebutuhan memori kecil dan waktu kompilasi
singkat. Kerugiannya, optimasi terbatas dan dukungan bahasa menjadi
kaku. Misalnya, pemanggilan fungsi sebelum deklarasinya sulit
ditangani, dan analisis lintas prosedur tidak mungkin dilakukan.

\subsection{Kompilator Banyak Lintasan}

Kompilator banyak lintasan (\eng{multi-pass compiler}) membaca
representasi program beberapa kali. Lintasan pertama dapat membangun
tabel simbol dan AST; lintasan berikutnya melakukan pemeriksaan tipe,
optimasi, dan pembangkitan kode. Pendekatan ini memungkinkan analisis
yang lebih dalam dan optimasi yang lebih baik.

\begin{figure}[H]
\centering
\begin{tikzpicture}[node distance=5mm and 8mm]
  \node[phase, text width=2.6cm] (a1) {Sumber};
  \node[phase, text width=2.6cm, right=of a1] (a2) {Kode target};
  \draw[arrow] (a1) -- node[above, font=\small] {satu lintasan} (a2);

  \node[phase, text width=2.2cm, below=of a1] (b1) {Sumber};
  \node[phase, text width=2.2cm, right=of b1] (b2) {AST};
  \node[phase, text width=2.2cm, right=of b2] (b3) {IR};
  \node[phase, text width=2.2cm, right=of b3] (b4) {Target};
  \draw[arrow] (b1) -- (b2);
  \draw[arrow] (b2) -- (b3);
  \draw[arrow] (b3) -- (b4);
\end{tikzpicture}
\caption{Perbandingan alur satu lintasan dan banyak lintasan}
\label{fig:ch03-pass}
\end{figure}

\subsection{Organisasi Front-End dan Back-End}

Pemisahan front-end dan back-end memungkinkan satu front-end dipakai
untuk beberapa target. Front-end memahami bahasa sumber dan
menghasilkan IR; back-end memahami IR dan menghasilkan kode untuk
arsitektur tertentu. Dengan desain ini, penambahan target baru tidak
memerlukan perubahan front-end, dan penambahan bahasa baru tidak
memerlukan perubahan back-end.

\begin{table}[H]
\centering
\caption{Trade-off one-pass dan multi-pass}
\label{tab:ch03-pass}
\begin{tabularx}{\textwidth}{@{}lXX@{}}
\toprule
\textbf{Aspek} & \textbf{One-pass} & \textbf{Multi-pass} \\
\midrule
Memori & Kecil & Lebih besar \\
Kecepatan kompilasi & Cepat & Lebih lambat \\
Kedalaman analisis & Terbatas & Mendalam \\
Optimasi & Minimal & Luas \\
Fleksibilitas bahasa & Rendah & Tinggi \\
Contoh & Pascal klasik & GCC, Clang, LLVM \\
\bottomrule
\end{tabularx}
\end{table}

\begin{warningbox}
One-pass bukan berarti sederhana. Banyak kompilator satu lintasan
tetap memerlukan tabel simbol dan mekanisme \eng{backpatching} untuk
menangani alamat yang belum diketahui. Trade-off utama tetap antara
memori, kecepatan, dan kualitas kode.
\end{warningbox}
